% RESEARCH_DRAFT_SHA256: 74f86a710b01cc8176663c6232476f240c21db50f82efe0744273ca0e0b672f6
% Layout and source-reference wrappers expanded for the web renderer.
\documentclass[11pt]{article}
\usepackage[margin=25mm]{geometry}
\usepackage{amsthm}
\usepackage{fontspec,unicode-math}
\setmainfont{Latin Modern Roman}
\setmathfont{Latin Modern Math}
\usepackage{amsmath,amssymb,mathtools}
\usepackage{xurl,hyperref}
\hypersetup{colorlinks=true,urlcolor=blue,linkcolor=blue}
\setcounter{secnumdepth}{0}
\setlength{\parindent}{0pt}
\setlength{\parskip}{5pt}
\setlength{\emergencystretch}{3em}
\newcommand{\sref}[2]{\hyperlink{source:#1:#2}{[S#2]}}
\newcommand{\cataloguescope}[1]{\paragraph{Source question.}#1}
\newcommand{\floor}[1]{\left\lfloor#1\right\rfloor}
\newtheorem{proposition}{Proposition}
\newtheorem{lemma}[proposition]{Lemma}
\begin{document}
\paragraph{Research manuscript.} OpenAI \texttt{gpt-6-astra}, reasoning \texttt{ultra}; 11 October 2026.
\paragraph{Preserved source body.} The complete catalogue mathematical text and sources follow with layout and source links expanded, including source qualifications and its historical preparation notes. New research begins at the section headed Research attempt.

\section{Multicolored Sunflower Conjecture over Ternary Vector Spaces}\label{statement}

\subsection{Definitions and mathematical statement}

\paragraph{Shared notation.}Standard number systems are denoted by $\mathbb N,\mathbb Z,\mathbb Q,\mathbb R,\mathbb C$; index ranges, computational models and specialized notation follow the cited problem source and the definitions below.


$\mathbb Z_3^n$ is the ternary vector space with coordinatewise addition modulo three. Every supplied triple sums to zero; the proposed conclusion mixes one coordinate from each of three selected triples.
\paragraph{Statement recorded in the supplied source.}
Does there exist $c>0$ such that every $S\subseteq(\mathbb Z_3^n)^3$ with $|S|\geq3^{n(1-c)}$ and $x+y+z=0$ for every $(x,y,z)\in S$ contains three elements $(x_i,y_i,z_i)$ satisfying $x_1+y_2+z_3=0$? \href{https://doi.org/10.4171/owr/2012/54}{[S2]}
\paragraph{Scope qualification.}
The supplied question does not explicitly say whether the three selected elements must be distinct; this boundary condition is not silently filled in.
\paragraph{Source status.}
The supplied catalogue labels this record open. Its source-status metadata was checked against the supplied snapshot on 2026-10-11. This import does not independently establish current open status or mathematical correctness. The catalogue literature assessment is dated 2026-08-21.
\subsection{Short English statement}
Find a cross-summing triple among three-coloured zero-sum triples once their set exceeds the stated exponential threshold.
\subsection{Sources}
\begin{itemize}
\item[S1] Ulam AI and the UnsolvedMath Contributors, UnsolvedMath: A Curated Collection of Open Mathematics Problems, catalogue entry. Catalogue curation and metadata: CC BY 4.0. The supplied source question is retained; definitions and the English summary are editorial additions.
\url{https://huggingface.co/datasets/ulamai/UnsolvedMath}.
\item[S2] Complexity Theory (2013), Oberwolfach Reports, DOI 10.4171/owr/2012/54, pp. 3267--3304
\url{https://doi.org/10.4171/owr/2012/54}.
\item[S3] sunflower and matrix multiplication paper
\url{https://people.csail.mit.edu/virgi/6.890/papers/sunflowersmult.pdf}.
\textit{Underlying Oberwolfach Reports material retains its source-specific licence, including CC BY-SA 4.0 where applicable; see THIRD\_PARTY\_NOTICES.md.}
\end{itemize}
\paragraph{Preparation and review.}
Original question attributed to the cited source and UnsolvedMath contributors. Codex (GPT-6; exact serving revision not exposed) assisted with definitions, English summaries, source inspection and formatting. No new human mathematical authorship or independent proof verification is asserted. Catalogue identity checked against the complete supplied registry and existing repository statements; known flagged repetitions and obvious identical conjectures excluded. Full mathematical-equivalence/literature review has not been performed. Statement preparation and final review followed contribution prompts C and D; automated checks verify structure and identity, not open status.
\label{end:statement}

\section{Research attempt}
\subsection{Literature, source context and scope}
The full original catalogue body and source qualifications are preserved above. The relevant primary workshop passages were read in the cached publisher report. Historical claims and subsequent attributed results are distinguished from the mathematical work below. A restricted example, standard lemma or known theorem is not counted as new progress on the entire original target. Missing assumptions and repeated or rhetorical extracts are not silently turned into different solved problems.

\subsection{Resolve the interpretation before applying the polynomial method}
The printed Conjecture 7 does not say whether its three selected elements must be distinct; the earlier uniquely solvable puzzle condition instead says that the three indices are not all equal. These readings differ. If arbitrary repetitions including three copies of one element are allowed, any nonempty $S$ works because its elements already satisfy $x+y+z=0$. If three pairwise distinct elements are required, the family
\[
S=\{(x,-x,0):x\in\mathbb F_3^n\}
\]
is a counterexample: $x_1+y_2+z_3=x_1-x_2=0$ forces the first two selected elements to coincide. Its size is $3^n$, exceeding every proposed $3^{(1-c)n}$ with $c>0$, $n\ge1$.
The standard tri-colored sum-free interpretation requires an off-diagonal choice of indices, meaning not all three equal. For this reading a later polynomial-method theorem applies. We give a complete, deliberately nonoptimal rederivation of that known bound, not a new discovery or a silent repair of the original quantifiers.
\begin{proposition}
Let $n\ge1$ and let $(a_i,b_i,c_i)\in(\mathbb F_3^n)^3$, $1\le i\le m$, satisfy
\[
a_i+b_j+c_k=0\quad\Longleftrightarrow\quad i=j=k.
\]
Then $m\le\theta^n$ where $\theta=(7/4)2^{2/3}<14/5<3$. In particular, the not-all-equal interpretation holds with $c=1-\log_3(14/5)>0$.
\end{proposition}
\begin{proof}
A slice of a three-index array is an expression $f(i)g(j,k)$, or one of the other two analogous forms. First we prove that an $m$ by $m$ by $m$ array equal to 1 when $i=j=k$ and zero otherwise requires at least $m$ slices over any field. Suppose it were expressed using $r$ slices separated in $i$, $s$ in $j$, and $t$ in $k$, with $r+s+t<m$. Let $U$ be the subspace of vectors $u\in\mathbb F^m$ annihilating the $r$ one-variable functions in the first type. Its dimension $d$ is at least $m-r$. Select $d$ coordinate functionals whose restrictions form a basis of $U^*$; prescribing all these coordinates to be 1 produces a vector $u\in U$ with at least $d$ nonzero coordinates. Contract the array against $u$ in its first index. The diagonal description gives a matrix with diagonal $u$, of rank at least $d$. The first $r$ slices vanish, while each of the remaining $s+t$ slices becomes a rank-at-most-one matrix. Thus its rank is at most $s+t<m-r\le d$, a contradiction. The $m$ single-coordinate slices show that the required number is exactly $m$.
Over $\mathbb F_3$, form the polynomial array
\[
T(i,j,k)=\prod_{\ell=1}^n\bigl(1-(a_{i\ell}+b_{j\ell}+c_{k\ell})^2\bigr).
\]
A factor is 1 exactly when the corresponding sum is zero and is 0 otherwise. The hypothesis therefore makes $T$ the diagonal array just analyzed. Its expansion has total degree at most $2n$, and each individual coordinate variable has exponent at most 2. In every monomial one of the three variable groups has degree at most $\lfloor2n/3\rfloor$. Assign each monomial to one such group, with a fixed tie rule, and collect monomials having the same factor in that group's variables. Each collection is one slice. Consequently
\[
m\le3D_n,\qquad D_n=\bigl|\{\alpha\in\{0,1,2\}^n:\sum_\ell\alpha_\ell\le2n/3\}\bigr|.
\]
For $0<t<1$, each counted exponent has $t^{\sum\alpha_\ell}\ge t^{2n/3}$. Summing and then including all exponents gives
\[
D_n\le t^{-2n/3}(1+t+t^2)^n.
\]
Taking $t=1/2$ yields $m\le3\theta^n$. To remove the factor 3, take the $r$-fold Cartesian product of the indexed family. A mixed sum vanishes precisely when the three index tuples agree in every component, so the same bound in dimension $nr$ gives $m^r\le3\theta^{nr}$. Taking $r$th roots and letting $r$ tend to infinity gives $m\le\theta^n$. Finally, $\theta^3=343/16<(14/5)^3<27$. A family of size at least $3^{(1-c)n}=(14/5)^n$ cannot be tri-colored sum-free, proving the last assertion.
\end{proof}
\subsection{Verification and remaining gaps}
The pairwise-distinct counterexample was enumerated for $n=1,2$. Exponent counts $D_n$ were checked by finite enumeration for $1\le n\le8$, and the strict cubic comparison was checked with exact rational arithmetic. The diagonal-slice argument was audited over a finite field: choosing pivot coordinates equal to 1 avoids an invalid appeal to a generic vector over an infinite field. The resulting exponential bound is a known consequence of the 2016 polynomial/slice-rank methods, and later work gives sharper constants. We claim no new progress beyond that literature. The original ambiguous phrase must still be interpreted explicitly; the three readings above are not interchangeable, and the catalogue status is not edited.

\section{References and source attribution}
Primary workshop report: \url{https://doi.org/10.4171/owr/2012/54}. Relevant passages read: 2012/54 pp. 3291--3293, uniquely solvable puzzles, the not-all-equal index convention, and Conjecture 7’s wording.

\url{https://arxiv.org/abs/1605.06702}. Read primary abstract of Blasiak, Church, Cohn, Grochow, Naslund, Sawin and Umans on tri-colored sum-free bounds and the cap-set method; the proof here is a self-contained rederivation of known methodology.

\url{https://arxiv.org/abs/1607.00047}. Read primary abstract of Kleinberg, Sawin and Speyer giving the standard iff-i=j=k definition and exponential bounds with sharper asymptotics.

All retrievals were checked on 11 October 2026. Original question attribution remains in the preserved Sources section. External theorems are attributed and not claimed as new discoveries or fully verified merely by citation.


\paragraph{Integration review.} Codex (GPT-6; exact serving revision not exposed) reviewed the written arguments and expanded layout and source-reference wrappers. Repository metadata and the production web renderer were checked using Python 3.13 and JavaScript with MathJax 3.2.2. This model-assisted review is not independent human proof verification; no PDF compilation or formal proof checking is asserted.

\section{Completion Estimate}
\noindent\textbf{COMPLETION: 0\%}
No substantiated new progress toward the entire original target is claimed. The worked results and failed route are delimited above; a known result, illustrative model, or auxiliary obstruction is not counted as original-target completion.

\section{Final status and attribution}
\textbf{Attempt status: unresolved.} This describes this research attempt, not a new ruling on the source catalogue's status or on external literature claims.
The mathematics and this self-audit were generated by OpenAI \texttt{gpt-6-astra}, reasoning \texttt{ultra}, on 11 October 2026. Codex (GPT-6; exact serving revision not exposed) orchestrated the batch, reviewed the written mathematical arguments, and contributed formatting/repository review. No human mathematical authorship, independent human proof checking, or formal verification is claimed. The historical preparation paragraphs in the reproduced catalogue body concern that body only.

Contribution prompt A was applied to the complete definition and supplied source record, followed by the adversarial checks of prompt B. Local reproduction files are \texttt{batch019/checks.py}, \texttt{computations.json}, and the per-ID audit JSON. They were not public links at preparation time. Numerical checks are finite and do not prove unrestricted results. No LaTeX engine was found locally; compilation was not performed by this research agent.

\end{document}
